\documentclass[tikz,border=6pt]{standalone}
\usepackage{tikz}
\usepackage{tikz-3dplot}

\begin{document}
\tdplotsetmaincoords{65}{115}
\begin{tikzpicture}[
  tdplot_main_coords,
  scale=1.2,
  line cap=round,
  line join=round,
  point/.style={circle,fill=black,inner sep=1.2pt},
  label/.style={font=\scriptsize,inner sep=1pt}
]
  % 3D coordinates:
  % a: x-axis at y=0,z=0
  % b: y-axis shifted to x=0,z=1
  % MN is perpendicular to both a and b, so MN is their common perpendicular.
  \coordinate (A1) at (-1.8,0,0);
  \coordinate (A2) at (1.8,0,0);
  \coordinate (B1) at (0,-1.3,1);
  \coordinate (B2) at (0,1.3,1);
  \coordinate (M) at (0,0,0);
  \coordinate (N) at (0,0,1);

  \draw[blue!75!black,line width=1.25pt] (A1) -- (A2) node[label,right] {$a$};
  \draw[orange!90!black,line width=1.25pt] (B1) -- (B2) node[label,above] {$b$};
  \draw[teal!70!black,line width=1.2pt] (M) -- (N) node[label,midway,left] {$d$};

  \fill[point] (M) node[label,below] {$M$};
  \fill[point] (N) node[label,above] {$N$};

  % Right-angle markers in 3D coordinate planes.
  \draw[black!65,line width=0.6pt]
    (0.28,0,0) -- (0.28,0,0.28) -- (0,0,0.28);
  \draw[black!65,line width=0.6pt]
    (0,0.28,1) -- (0,0.28,0.72) -- (0,0,0.72);

\end{tikzpicture}
\end{document}
